% \iffalse meta-comment
%
% hit-report-keylist.dtx 是开题中期报告类的键清单：类选项、信息字段、常量
%
% \fi
%
% \section{开题中期报告类的键清单（keylist）}
%
% \changes{v3.2a}{0000/00/00}{元信息字段的声明单独成文件。它说的是「有哪些字段」，
%   与 \file{src/setup/} 说的「字段取什么值」是一对，都是数据；
%   字段怎么存、旧名怎么兼容那套机制在 \file{src/utils/hit-meta.dtx}}
% 开题中期报告用到的元信息字段：字段名、旧名映射、\cs{hitsetup} 各族的路由，
% 以及摘要环境。都是清单，不含机制。
%
% ^^A ======================================================================
% ^^A Module report-keylist: 元信息字段（\hitsetup 接受的 key）（hit-report）
% ^^A ======================================================================
%    \begin{macrocode}
%<*report-keylist>
%% 中文封面

\hit@define@title{zh}
\hit@define@term{report-type-zh}
\hit@define@term{specialty-zh}
\hit@define@term{author-zh}
\hit@define@term{student-id}
\hit@define@term{class-id}
\hit@define@term{affiliation-zh}
\hit@define@term{supervisor-zh}
\hit@define@term{date-zh}
\hit@define@term{enroll-date-zh}

%    \end{macrocode}
%
% \changes{v3.2a}{0000/00/00}{关键词的逗号拆分与 \cs{hitsetup} 的入口挪到
%   \file{src/utils/hit-meta.dtx}}
% 关键词怎么拆、\cs{hitsetup} 的占位定义都见 \file{src/utils/hit-meta.dtx}。
%
% \changes{v3.2a}{0000/00/00}{\cs{hitsetup} 改用 l3keys 做入口，加一层按职责
%   划分的子族路由，与 \file{hit-thesis.cls} 一致。子族暂时只转发，键仍由
%   \texttt{hit/class} 族处理}
% 入口换成 l3keys，根族路由到子族，子族一律把键转回
% \texttt{hit} 族。转发时值要重新加花括号，否则含逗号的值会被二次拆分成多个键。
%
% \changes{v3.2a}{0000/00/00}{元信息、可本地化字符串与参考文献选项正式注册进
%   对应的族，不再依赖子族兜底转发}
% 元信息字段与可本地化字符串。做法与 \file{hit-thesis.cls} 相同：
% 前者的处理器直接调用同名命令，后者的键名由宏名去掉 \texttt{hit@} 前缀、
% 其余 \texttt{@} 换成连字符得来。
%    \begin{macrocode}
\hit@declare@info@field
  {
    title-zh , affiliation-zh , author-zh , class-id , date-zh , enroll-date-zh , student-id ,
    specialty-zh , supervisor-zh , report-type-zh
  }
\hit@declare@constant@as
  {
    equationname , listfigureename , listtableename , listequationename ,
    listequationname , cabstractcname , cabstractename , eabstractcname , eabstractename
  }
\hit@declare@constant@ctex
  {
    chapter-name = chapter/name , appendixname , contentsname , listfigurename ,
    listtablename , figurename , tablename , bibname , indexname
  }
\hit@declare@constant
  {
    campus-full , algorithm-name-zh , algorithm-name-en ,
    report-header , report-front-header ,
    bibliography-heading-zh ,
    authorization-usage-heading , author-signature-field-label , date-blank ,
    date-field-label-front , harbin-bachelor-school-bottom-mark , harbin-school-bottom-mark ,
    shenzhen-doctor-midterm-note , shenzhen-master-school-bottom-mark ,
    supervisor-signature-field-label , bachelor-affiliation-field-label ,
    bachelor-affiliation-field-label-shenzhen , bachelor-affiliation-field-label-weihai ,
    bachelor-class-field-label , bachelor-date-field-label ,
    bachelor-check-date-field-label , bachelor-major-field-label ,
    bachelor-student-id-field-label , bachelor-student-field-label ,
    bachelor-supervisor-field-label , bachelor-document-type , bachelor-degree-level ,
    bachelor-teacher-comment-field-label , bachelor-teacher-signature-field-label ,
    affiliation-field-label-zh , associate-supervisor-field-label-zh , author-field-label-zh ,
    co-supervisor-field-label-zh , date-field-label-zh , degree-of-discipline-format-zh ,
    degree-field-label-zh , school-zh , school-field-label-zh , specialty-field-label-zh ,
    supervisor-field-label-zh , thesis-midterm-check , document-type ,
    thesis-title-prefix , degree-level-zh , figure-prefix-en , table-prefix-en ,
    eqdenote-label-zh , eqdenote-label-en , eqdenote-dash , graduate-affiliation-field-label
    , graduate-affiliation-major-field-label , graduate-major-field-label ,
    graduate-date-field-label , graduate-enroll-date-field-label ,
    graduate-student-id-field-label , graduate-student-field-label ,
    graduate-supervisor-field-label , graduate-thesis-field-label , hi ,
    classified-index-international-field-label-zh , classified-index-national-field-label-zh ,
    school-id , school-id-field-label , secret-level-field-label , campus ,
    stage-document-type , stage-midterm , stage-opening , title-separator-zh
  }
%    \end{macrocode}
%
% 参考文献的三个选项，改用连字符命名。
%    \begin{macrocode}
\keys_define:nn { hit / bib }
  {
    label-align .code:n = { \str_set:Nn \g_bib_labelalign_str {#1} } ,
    author-case .code:n =
      { \str_if_eq:nnT {#1} { full } { \bool_set_true:N \g_bib_authorupper_bool } } ,
    max-author .code:n =
      {
        \int_set:Nn \g_bib_maxauthor_int {#1}
        \int_if_zero:nT { \g_bib_maxauthor_int } { \int_gdecr:N \g_bib_maxauthor_int }
      } ,
  }
\keys_define:nn { hit / style }
  {
    ragged-bottom .choice: ,
    ragged-bottom / true  .code:n = { \bool_gset_true:N \g__hit_raggedbottom_bool } ,
    ragged-bottom / false .code:n = { \bool_gset_false:N \g__hit_raggedbottom_bool } ,
    ragged-bottom .default:n = true ,
  }
%    \end{macrocode}
%
% \changes{v3.2a}{0000/00/00}{加旧名到新名的映射表与废弃警告}
% 兼容期的处理，与 \file{hit-thesis.cls} 一致：旧名照常生效，只在日志里提示。
%    \begin{macrocode}
\prop_gset_from_keyval:Nn \g__hit_renamed_option_prop
  {
    raggedbottom   = { style = { ragged-bottom = ... } } ,
    biblabelalign  = { bib = { label-align = ... } } ,
    bibauthorupper = { bib = { author-case = ... } } ,
    bibmaxauthor   = { bib = { max-author = ... } } ,
  }
%    \end{macrocode}
%
% \changes{v3.2a}{0000/00/00}{元信息字段的旧名走废弃期：仍然可用，用一次提示一次}
% 旧名到新名的对照。旧名在 v3.1e 就是公开接口，写在用户的 \file{cover.tex} 里，
% 所以保留到 v4.1。
%    \begin{macrocode}
\hit@define@legacy@terms
  {
    ctype=report-type-zh ,
    csubject=specialty-zh ,
    cauthor=author-zh ,
    cstudentid=student-id ,
    cclassid=class-id ,
    caffil=affiliation-zh ,
    csupervisor=supervisor-zh ,
    cdate=date-zh ,
    cenrolldate=enroll-date-zh
  }
\hit@define@legacy@title{ctitlecover}{zh}{first}
\hit@define@legacy@title{ctitleone}{zh}{first}
\hit@define@legacy@title{ctitletwo}{zh}{second}
\prop_gset_from_keyval:Nn \g__hit_renamed_load_option_prop
  {
    font     = { font-en = ... } ,
    cjk-font = { font-zh = ... } ,
    type     = { level } ,
  }
\__hit_check_class_options:
%    \end{macrocode}
%
% 根族路由。\option{toc} 既是族名又是老的布尔选项名，所以它单独处理：
% 值是 \texttt{true}/\texttt{false} 时当老写法，其余当子族的键值表。
%    \begin{macrocode}
\keys_define:nn { hit }
  {
    info  .meta:nn = { hit / info  } {#1} ,
    style .meta:nn = { hit / style } {#1} ,
    bib   .meta:nn = { hit / bib   } {#1} ,
    const .meta:nn = { hit / const } {#1} ,
    structure .meta:nn = { hit / structure } {#1} ,
    theorem .code:n = { \hit@theorem@style {#1} } ,
    theorem-list .code:n = { \hit@theorem@list {#1} } ,
    theorem-star .code:n = { \hit@theorem@declare@star@list {#1} } ,
    toc .code:n =
      {
        \str_case:nnF {#1}
          {
            { true }  { \bool_gset_true:N \g__hit_toc_bool }
            { false } { \bool_gset_false:N \g__hit_toc_bool }
          }
          { \keys_set:nn { hit / toc } {#1} }
      } ,
    toc .default:n = true ,
    unknown .code:n = { \__hit_legacy_key:n {#1} } ,
  }
\RenewDocumentCommand \hitsetup { O{} +m }
  { \hit@if@condition:nT {#1} { \keys_set:nn { hit } {#2} } }
%</report-keylist>
%    \end{macrocode}
%
%
% \Finale